\documentclass[11pt]{article}
\usepackage{amsmath,amssymb,amsthm}
\usepackage{mathtools}
\usepackage{tikz}
\usepackage{enumitem}
\usepackage[margin=1in]{geometry}

\newtheorem{theorem}{Theorem}
\newtheorem*{theorem*}{Theorem}
\newtheorem{lemma}{Lemma}
\newtheorem{definition}{Definition}
\newtheorem*{definition*}{Definition}
\newtheorem{proposition}{Proposition}
\newtheorem{assumption}{Assumption}
\newtheorem{remark}{Remark}



\title{Math 106 -- Notes \\ Week 6: February 10/12, 2026}
\author{Ethan Levien}
\date{}

\begin{document}
\maketitle


\section*{Stochastic differential equations: Existence and uniqueness}
\begin{theorem}
Assume the coefficients 
\[
b:\mathbb{R}^n\times[0,T]\to\mathbb{R}^n,
\qquad
\sigma:\mathbb{R}^n\times[0,T]\to\mathbb{R}^{n\times m},
\]
satisfy the global Lipschitz condition
\begin{equation}
\|b(x,t)-b(y,t)\| + \|\sigma(x,t)-\sigma(y,t)\|
\;\le\; K\,\|x-y\|
\end{equation}
and the linear growth condition
\begin{equation}
\|b(x,t)\|^2 + \|\sigma(x,t)\|^2 
\;\le\; K^2\,(1+\|x\|^2)
\end{equation}
for all $x,y\in\mathbb{R}^n$ and $t\in[0,T]$.  
Here $\|\sigma\|$ denotes the Frobenius norm:
\[
|b|^2 := \sum_{i=1}^n b_i^2,
\qquad
\|\sigma\|^2 := \sum_{i=1}^n \sum_{j=1}^m \sigma_{ij}^2.
\]
Then the equation
\begin{equation}
X_t = X_0 + \int_0^t b(X_s,s)\,ds + \int_0^t \sigma(X_s,s)\,dW_s
\end{equation}
has a unique solution such that each component of $X_t$ belongs to $\mathcal{V}$.  
Moreover, the process $X_t$ is almost surely continuous in $t$.
\end{theorem}

We will need to recall the following result. 

\begin{theorem}[Gr\"onwall's inequality]
Let $u:[0,T]\to[0,\infty)$ be a continuous function satisfying
\begin{equation}
u(t)
\;\le\;
a
+
\int_0^t b(s)\,u(s)\,ds,
\end{equation}
for all $t\in[0,T]$, where $a\ge 0$ is a constant and $b:[0,T]\to[0,\infty)$ is integrable.  
Then
\begin{equation}
u(t)
\;\le\;
a\,\exp\!\left( \int_0^t b(s)\,ds \right),
\qquad 0 \le t \le T.
\end{equation}
In particular, if $a=0$ then $u(t)\equiv 0$ on $[0,T]$.
\end{theorem}

\begin{proof}[Proof of uniqueness] Following the book, we will only do the one-dimensional case and neglect explicit time dependence in $b$ and $\sigma$. I will show uniqueness, showing existence is more complicated. We need to show that if $X_t$ and $\hat{X}_t$ are solutions of the SDE with the same initial value, then they are almost surely equal. Using $(a+b)^2 \le 2a^2 + 2b^2$ and taking expectations, we get
\begin{equation}
\mathbb{E}\big[(X_t - \hat{X}_t)^2\big]
\le 
2\,\mathbb{E}\left[\left( \int_0^t \big(b(X_s) - b(\hat{X}_s)\big)\,ds\right)^2\right]
+
2\,\mathbb{E}\left[\left( \int_0^t \big(\sigma(X_s) - \sigma(\hat{X}_s)\big)\,dW_s\right)^2\right].
\end{equation}

For the drift term, apply Cauchy--Schwarz with the second function equal to $1$:
\begin{equation}
\left( \int_0^t \big(b(X_s) - b(\hat{X}_s)\big)\,ds \right)^2
\le 
t \int_0^t \big(b(X_s) - b(\hat{X}_s)\big)^2\,ds.
\end{equation}
Taking expectations and using the Lipschitz condition $|b(x)-b(y)| \le K|x-y|$,
\begin{equation}
\mathbb{E}\left[\left( \int_0^t \big(b(X_s) - b(\hat{X}_s)\big)\,ds \right)^2\right]
\le 
K^2 t \int_0^t \mathbb{E}\big[(X_s - \hat{X}_s)^2\big]\,ds.
\end{equation}

For the stochastic term, we use the It\^o isometry and  Lipschitz condition
\begin{align}
\mathbb{E}\left[\left( \int_0^t \big(\sigma(X_s) - \sigma(\hat{X}_s)\big)\,dW_s\right)^2\right]
&=
\mathbb{E}\left[ \int_0^t \big(\sigma(X_s) - \sigma(\hat{X}_s)\big)^2\,ds \right]\\
&\le
K^2 \int_0^t \mathbb{E}\big[(X_s - \hat{X}_s)^2\big]\,ds.
\end{align}

Combining these bounds, we obtain
\begin{equation}
\mathbb{E}\big[(X_t - \hat{X}_t)^2\big]
\le
2K^2 t \int_0^t \mathbb{E}\big[(X_s - \hat{X}_s)^2\big]\,ds
+
2K^2 \int_0^t \mathbb{E}\big[(X_s - \hat{X}_s)^2\big]\,ds,
\end{equation}
that is,
\begin{equation}
\mathbb{E}\big[(X_t - \hat{X}_t)^2\big]
\le
2K^2(t+1) \int_0^t \mathbb{E}\big[(X_s - \hat{X}_s)^2\big]\,ds.
\end{equation}

Define
\begin{equation}
u(t) = \mathbb{E}\big[(X_t - \hat{X}_t)^2\big].
\end{equation}
The last inequality shows that $u(t)$ satisfies
\begin{equation}
u(t) \le \int_0^t 2K^2(t+1)\,u(s)\,ds.
\end{equation}
By Gr\"onwall's inequality with zero initial term, it follows that $u(t) = 0$ for all $t \in [0,T]$. Hence
\begin{equation}
\mathbb{E}\big[(X_t - \hat{X}_t)^2\big] = 0,
\end{equation}
which implies $X_t = \hat{X}_t$ almost surely for each $t$, and in particular for all rational $t$. By continuity of the sample paths, $X_t = \hat{X}_t$ almost surely for all $t\in[0,T]$. 
\end{proof}




\section*{Numerical schemes}

\subsection*{Brief review of numerical solutions to ODE}

Consider the ODE
\begin{equation}
\frac{d}{dt} x_t = b(x_t).
\end{equation}
We formulate numerical algorithms by taking the integral form 
\begin{equation}
x_{t+dt} = x_t + \int_t^{t+dt} b(x_s)\,ds
\end{equation}
and approximating the integral: 
\begin{equation}
 \int_t^{t+dt} b(x_s)\,ds 
 =
 \sum_{j=1}^n b(x_{t+h_j})\, w_j,
 \qquad
 h_j \in [0,dt).
\end{equation}
Here the nodes $h_j$ and weights $w_j$ define a quadrature rule on $[0,dt]$.

On the other hand, we have the Taylor expansion 
\begin{equation}
x_{t + dt} 
=
x_t 
+
\sum_{i=1}^{\infty}
\frac{(dt)^i}{i!}
\frac{d^i}{dt^i}x_t.
\end{equation}
This can be written in operator form. Consider 
\begin{equation}
\frac{d}{dt} f(x_t) 
= f'(x_t)b(x_t) 
= ({\mathcal L} f)(x_t),
\end{equation}
where ${\mathcal L}$ is the generator (Koopman operator)
\begin{equation}
{\mathcal L} f(x) = b(x) f'(x).
\end{equation}
Hence time derivatives of functions of $x_t$ are obtained by iterative application of ${\mathcal L}$:
\begin{equation}
x_{t + dt} 
=
x_t 
+
\sum_{i=1}^{\infty}
\frac{(dt)^i}{i!}
{\mathcal L}^i f(x_t),
\qquad
f(x)=x.
\end{equation}
Equivalently,
\begin{equation}
x_{t + dt} 
=
x_t 
+
\sum_{i=1}^{\infty}
\frac{(dt)^i}{i!}
{\mathcal L}^i x_t.
\end{equation}
Note that ${\mathcal L}$ acts on $x_t$ as the function $f(x)=x$ of the state variable $x$, not on the time variable $t$.

For a $k$th order method, the quadrature rule must reproduce the Taylor expansion up to order $k$, i.e.
\begin{equation}
\sum_{j=1}^n b(x_{t+h_j})\, w_j
=
\sum_{i=1}^{k}
\frac{(dt)^i}{i!}
{\mathcal L}^i x_t
+
\mathcal O(dt^{k+1}).
\end{equation}
Runge--Kutta methods are constructed precisely so that this matching holds without explicitly computing derivatives of $b$. A key point is that the $h_j$ depend on $x$. 

\medskip

The two most important examples are:

\begin{itemize}

\item Euler's method: This corresponds to the left endpoint rule and is first order,
\begin{equation}
x_{n+1} = x_n + dt\, b(x_n).
\end{equation}

\item Runge--Kutta 4 (fourth order): This method evaluates $b$ at several intermediate points in order to approximate the integral
\(
\int_t^{t+dt} b(x_s)\,ds
\)
with fourth-order accuracy. Writing
\begin{align}
k_1 &= b(x_n), \\
k_2 &= b\!\left(x_n + \frac{dt}{2}k_1\right), \\
k_3 &= b\!\left(x_n + \frac{dt}{2}k_2\right), \\
k_4 &= b(x_n + dt\, k_3),
\end{align}
the update is
\begin{equation}
x_{n+1}
=
x_n
+
\frac{dt}{6}
\left(
k_1 + 2k_2 + 2k_3 + k_4
\right).
\end{equation}

Each stage $k_j$ approximates the vector field along an approximate trajectory inside the time interval $[t,t+dt]$. The particular weights
\(
\frac{1}{6}(1,2,2,1)
\)
are chosen so that, when the method is expanded in powers of $dt$, it agrees with the Taylor expansion of the true solution up to terms of order $dt^4$. Thus RK4 achieves fourth-order accuracy using only evaluations of $b$, without requiring $b'$ or higher derivatives.

\end{itemize}

\subsection*{The It\^o-Taylor expansion}

For the SDE
\begin{equation}
dX_t = b(X_t)\,dt + \sigma(X_t)\,dW_t,
\end{equation}
the structure is formally the same as in the deterministic case, but now the flow map is random.  Writing the integral form,
\begin{equation}
X_{t+dt}
=
X_t
+
\int_t^{t+dt} b(X_s)\,ds
+
\int_t^{t+dt} \sigma(X_s)\,dW_s,
\end{equation}
we again seek to approximate the integrals.  The difference is that the (1) second integral is stochastic and must be simulated and (2) we need to apply It\^o's formula to obtain the Taylor expansion, which will now be called the It\^o-Taylor expansion. 

In particular, for a suitable function $f$, 
 It\^{o}'s formula gives
\begin{equation}
df(X_t)
=
(\mathcal L_1 f)(X_t)
+
(\mathcal L_2 f) (X_2)
\end{equation}
where we now have two operators
\begin{align}
\mathcal L_1 f(x)
&=
b(x) f'(x)
+
\frac12 \sigma(x)^2 f''(x)\\
\mathcal L_2 f(x) &= \sigma(x)f'(x)
\end{align}
and It\^{o}'s formula gives
\begin{equation}
df(X_t)
=
(\mathcal L_1 f)(X_t)\,dt
+
(\mathcal L_2 f)(X_t)dW_t.
\end{equation}
Iterating this identity produces the It\^{o}--Taylor expansion.  The key structural difference from the ODE case is the scaling
\begin{equation}
\Delta W := W_{t+dt}-W_t \sim {\rm Normal}(0,dt),
\qquad
\Delta W = O(dt^{1/2}),
\qquad
(\Delta W)^2 = dt + O(dt^{3/2}),
\end{equation}
so the expansion is graded in powers of $dt^{1/2}$ rather than $dt$.


Integrating from $t$ to $t+dt$ gives
\begin{equation}\label{eq:first_step}
f(X_{t+dt})
=
f(X_t)
+
\int_t^{t+dt} (\mathcal L_1 f)(X_s)\,ds
+
\int_t^{t+dt} (\mathcal L_2 f)(X_s)\,dW_s.
\end{equation}
Just as with Taylor expansion, we now expand the integrands 
$(\mathcal L_i f)(X_s)$ for $i\in\{1,2\}$ around time $t$ by applying It\^{o}'s formula again.  
Let $g(x) := (\mathcal L_i f)(x)$. Then
\begin{equation}
d g(X_u)
=
(\mathcal L_1 g)(X_u)\,du
+
(\mathcal L_2 g)(X_u)\,dW_u.
\end{equation}
Integrating from $u=t$ to $u=s$ gives
\begin{equation}
(\mathcal L_i f)(X_s)
=
(\mathcal L_i f)(X_t)
+
\int_t^s (\mathcal L_1 \mathcal L_i f)(X_u)\,du
+
\int_t^s (\mathcal L_2 \mathcal L_i f)(X_u)\,dW_u.
\end{equation}

Substituting this back into \eqref{eq:first_step} yields
\begin{align}
f(X_{t+dt})
&=
f(X_t)
+
(\mathcal L_1 f)(X_t)\,dt
+
(\mathcal L_2 f)(X_t)\,\Delta W \nonumber \\
&\quad
+
\int_t^{t+dt}
\int_t^{s}
(\mathcal L_1 \mathcal L_1 f)(X_u)\,du\,ds
+
\int_t^{t+dt}
\int_t^{s}
(\mathcal L_2 \mathcal L_1 f)(X_u)\,dW_u\,ds \nonumber \\
&\quad
+
\int_t^{t+dt}
\int_t^{s}
(\mathcal L_1 \mathcal L_2 f)(X_u)\,du\,dW_s
+
\int_t^{t+dt}
\int_t^{s}
(\mathcal L_2 \mathcal L_2 f)(X_u)\,dW_u\,dW_s.
\end{align}

We now recognize the appearance of iterated time and stochastic integrals.  
Freezing coefficients at $X_t$ (which is sufficient for deriving a truncated scheme) gives
\begin{align}
f(X_{t+dt})
&=
f(X_t)
+
(\mathcal L_1 f)(X_t)\,dt
+
(\mathcal L_2 f)(X_t)\,\Delta W \nonumber \\
&\quad
+
(\mathcal L_1^2 f)(X_t)\, I_{(1,1)}
+
(\mathcal L_2 \mathcal L_1 f)(X_t)\, I_{(2,1)} \nonumber \\
&\quad
+
(\mathcal L_1 \mathcal L_2 f)(X_t)\, I_{(1,2)}
+
(\mathcal L_2^2 f)(X_t)\, I_{(2,2)}
+
\mathcal O(dt^{3/2}),
\end{align}
where the iterated integrals are defined by
\begin{align}
I_{(1,1)} &= \int_t^{t+dt}\!\!\int_t^{s} du\,ds = \frac12 (dt)^2, \\
I_{(2,1)} &= \int_t^{t+dt}\!\!\int_t^{s} dW_u\,ds, \\
I_{(1,2)} &= \int_t^{t+dt}\!\!\int_t^{s} du\,dW_s, \\
I_{(2,2)} &= \int_t^{t+dt}\!\!\int_t^{s} dW_u\,dW_s
= \frac12\big((\Delta W)^2 - dt\big).
\end{align}

This is the It\^{o}--Taylor expansion up to strong order $1$.  
Truncating this expansion at different levels yields the Euler--Maruyama and Milstein schemes.



\subsection*{Euler--Maruyama and Milstein schemes}

Specializing the It\^{o}--Taylor expansion to $f(x)=x$, we have
\begin{equation}
\mathcal L_1 f(x)=b(x),
\qquad
\mathcal L_2 f(x)=\sigma(x),
\qquad
\mathcal L_2^2 f(x)=\sigma(x)\sigma'(x).
\end{equation}

Keeping terms in the expansion according to their order in $dt^{1/2}$ yields the following schemes:

\begin{itemize}

\item \textbf{Euler--Maruyama (strong order $1/2$):}

Truncate after the terms involving $dt$ and $\Delta W$,
\begin{equation}
X_{n+1}
=
X_n
+
(\mathcal L_1 f)(X_n)\,dt
+
(\mathcal L_2 f)(X_n)\,\Delta W_n,
\end{equation}
that is,
\begin{equation}
X_{n+1}
=
X_n
+
b(X_n)\,dt
+
\sigma(X_n)\,\Delta W_n,
\qquad
\Delta W_n \sim \mathrm{Normal}(0,dt).
\end{equation}

\item \textbf{Milstein (strong order $1$):}

Retain also the iterated integral $I_{(2,2)}$ term,
\begin{equation}
X_{n+1}
=
X_n
+
(\mathcal L_1 f)(X_n)\,dt
+
(\mathcal L_2 f)(X_n)\,\Delta W_n
+
\frac12 (\mathcal L_2^2 f)(X_n)\big((\Delta W_n)^2 - dt\big),
\end{equation}
that is,
\begin{equation}
X_{n+1}
=
X_n
+
b(X_n)\,dt
+
\sigma(X_n)\,\Delta W_n
+
\frac12 \sigma(X_n)\sigma'(X_n)
\big((\Delta W_n)^2 - dt\big).
\end{equation}

\end{itemize}
\noindent
\textbf{Strong convergence (mean square convergence).}
Let $X_t^{\delta t}$ denote the numerical solution of the SDE with (maximal) step size $\delta t$, and let $X_t$ denote the exact solution.  
We say that $\{X_t^{\delta t}\}$ converges strongly to $X_t$ with order $\alpha$ if
\begin{equation}
\max_{0 \le t \le T}
\mathbb E \big| X_t^{\delta t} - X_t \big|^2
\le
C (\delta t)^{2\alpha},
\end{equation}
where $C$ is a constant independent of $\delta t$.

Equivalently,
\begin{equation}
\max_{0 \le t \le T}
\left(
\mathbb E \big| X_t^{\delta t} - X_t \big|^2
\right)^{1/2}
\le
C (\delta t)^{\alpha}.
\end{equation}

Euler--Maruyama has strong order $\alpha = \tfrac12$, while the Milstein scheme has strong order $\alpha = 1$.

\medskip

\noindent
\textbf{Weak convergence (convergence with respect to expectation).}
We say that $\{X_t^{\delta t}\}$ converges weakly to $X_t$ with order $\beta$ if for every $f \in C_b^\infty(\mathbb R)$,
\begin{equation}
\max_{0 \le t \le T}
\left|
\mathbb E f(X_t^{\delta t})
-
\mathbb E f(X_t)
\right|
\le
C_f (\delta t)^{\beta},
\end{equation}
where $C_f$ is independent of $\delta t$ (but may depend on $f$).

Euler--Maruyama has weak order $\beta = 1$, even though its strong order is $\alpha = \tfrac12$.
\medskip

\noindent
\textbf{Interpretation.}
Strong order controls pathwise error and therefore requires matching terms in
the It\^{o}--Taylor expansion involving powers of $\Delta W$.
Weak order only requires matching expectations, so many stochastic terms vanish
after taking expectation.  As a result, weak schemes can achieve higher order
with less structure than strong schemes.

Thus the distinction is intrinsic to the stochastic setting: in deterministic
ODEs, strong and weak notions coincide, while for SDEs they differ because the
object being approximated is a probability law rather than a single trajectory.
\end{document}